\subsection{Измерительные приборы}
%%%%%%%%%%%%%%%%%%%%%%%%%%%%%%%%%%%%%%%%%%%%%%%%%%%%
%%%%%%%%%%%%%%%%%%%%%%%%%%%%%%%%%%%%%%%%%%%%%%%%%%%%
\z{К амперметру с собственным сопротивлением $r_\text{а}$ подключили шунт $R_\text{ш}$. Во сколько раз увеличился предел измерений?}
%
{$n = 1 + \frac{r_a}{r_{\text{Ш}}}$}

%%%%%%%%%%%%%%%%%%%%%%%%%%%%%%%%%%%%%%%%%%%%%%%%%%%%
\z{Какой шунт нужно подключить к амперметру, чтобы увеличить его предел измерений в $n$ раз? Сопротивление амперметра $r_\text{а}$.}
%
{$r_{\text{ш}} = \frac{r_{\text{А}}}{n-1}$}

%%%%%%%%%%%%%%%%%%%%%%%%%%%%%%%%%%%%%%%%%%%%%%%%%%%%
\z{К вольтметру с собственным сопротивлением $R_\text{в}$ подключили добавочное сопротивление $R_\text{д}$. Во сколько раз увеличился предел измерений?}
%
{$n = 1 + \frac{R_{\text{д}}}{R_{\text{В}}}$}

%%%%%%%%%%%%%%%%%%%%%%%%%%%%%%%%%%%%%%%%%%%%%%%%%%%%
\z{Какое добавочное сопротивление нужно подключить к вольтметру, чтобы увеличить его предел измерений в $n$ раз? Сопротивление вольтметра $R_\text{в}$.}
%
{$R_{\text{д}} = R_{\text{В}} (n-1)$}

%%%%%%%%%%%%%%%%%%%%%%%%%%%%%%%%%%%%%%%%%%%%%%%%%%%%
\z[{\begin{wrapfigure}[5]{r}{0.3\textwidth}%
		\vspace{-15pt}%
		\begin{flushright}%
 			\includegraphics[width=0.3\textwidth]{3.3.5}%
		\end{flushright}%
	\end{wrapfigure}%
	}]%
{В цепь включены последовательно три проводника сопротивлением $R_1 = 5\,\text{Ом}$, $R_2 = 6\,\text{Ом}$, $R_3 = 12\,\text{Ом}$. Какую силу тока показывает амперметр и какое напряжение между точками $A$ и $B$, если показания вольтметра $U = 1.2\,\text{В}$?
\vspace{3em}
}
%
{$I = 0.2\,\text{А};\quad V = 4.6\,\text{В}$}

%%%%%%%%%%%%%%%%%%%%%%%%%%%%%%%%%%%%%%%%%%%%%%%%%%%%
\z[{\begin{wrapfigure}[3]{r}{0.3\textwidth}%
		\vspace{-20pt}%
		\begin{flushright}%
 			\includegraphics[width=0.3\textwidth]{3.3.6}%
		\end{flushright}%
	\end{wrapfigure}%
	}]%
{В цепь включены два проводника $R_1 = 5\,\text{Ом}$, $R_2 = 10\,\text{Ом}$. Вольтметр $V_1$ показывает $12\,\text{В}$. Определить показания амперметра и вольтметра $V_2$.
\vspace{3em}
}
%
{$I = 0.8\,\text{А};\quad V_2 = 8\,\text{В}$}
%%%%%%%%%%%%%%%%%%%%%%%%%%%%%%%%%%%%%%%%%%%%%%%%%%%%
\z[{\begin{wrapfigure}[3]{r}{0.3\textwidth}%
		\vspace{-20pt}%
		\begin{flushright}%
 			\includegraphics[width=0.3\textwidth]{3.3.7}%
		\end{flushright}%
	\end{wrapfigure}%
	}]%
{Батарейка имеет напряжение $10\,\text{В}$. В схеме использованы одинаковые вольтметры, найти их показания.}
%
{$6\,\text{В},\ 4\,\text{В}$}
